\chapter{Regularity of the components}\label{ch:regular}
\chapterbrief
{the factor decomposition of Theorem~\ref{thm:factor}; $\Omega$ bounded; $u$ in a Sobolev space $H^m(\Omega)$ (square-integrable derivatives up to order $m$).}
{the decomposition is algebraic. To use the components numerically we need them to be as regular as $u$ itself, and to depend continuously on $u$.}
{completes the answer to Q2. It gives the exact condition (a conditioning number $\gamma_N>0$), shows that it can fail, and that the geometry of $\Omega$ against the characteristics decides it. Step~5 of the algorithm checks it.}

\section{The pieces can be rougher than the whole}

\begin{falsebox}{Loss of a derivative \led{C12-C1} \falsest}\label{cex:loss}
The statement ``if $u\in H^m(\Omega)$ then every component $u_j\in H^m(\Omega)$'' is false. Let $\Omega=(0,1)^2$ and
\[
f=\tau-\xi,\qquad g=\tau-\xi-1 .
\]
For every $m$ there is a smooth $u\in H^m(\Omega)$ with $f(\partial)g(\partial)u=0$ whose unique components $u_1\in\ker f$, $u_2\in\ker g$ are \emph{not} in $H^m$. Exactly one derivative is lost: $u_2=f(\partial)u\in H^{m-1}$.

\dimref{C12-C1}
\end{falsebox}

\begin{intuition}
In the coordinates $s=x+t$ along the characteristics, both factors are first-order ODEs along the same lines, with solutions $1$ and $e^{t}$. On a long chord these two functions are clearly different. Near a corner of the square the chords become arbitrarily short, and on a short chord $1$ and $e^{t}$ are almost collinear. Separating a function into these two shapes then amplifies it by a factor that blows up as the chord shrinks (Figure~\ref{fig:chords}).
\end{intuition}

\begin{figure}[!htb]
\centering
\begin{tikzpicture}[scale=2.6]
\draw[thick,cA] (0,0) rectangle (1,1);
\foreach \s in {0.12,0.3,0.5,0.7,0.88}{
  \draw[cB,thick] ({max(0,-\s+0.0)+0},{\s}) -- ({1-\s},1);}
\foreach \s in {0.12,0.3,0.5,0.7,0.88}{
  \draw[cC,thick] ({\s},0) -- (1,{1-\s});}
\node[font=\scriptsize,anchor=east] at (0,0.95) {short};
\node[font=\scriptsize,anchor=west] at (1,0.05) {short};
\node[font=\scriptsize] at (0.5,-0.12) {$t$};
\node[font=\scriptsize] at (-0.1,0.5) {$x$};
\end{tikzpicture}
\caption{The unit square and the characteristic chords of $\partial_t+\partial_x$. Chords near the corners $(0,1)$ and $(1,0)$ are arbitrarily short. On a short chord the two shapes of the components are nearly collinear, and the splitting is ill-conditioned.}\label{fig:chords}
\end{figure}

\section{Two components: the chord criterion}

For two transports in the same direction, the geometry decides everything.

\begin{theorem}{Chord criterion \led{C12-T3} \proved}\label{thm:chord}
Let $\Omega$ be bounded, open and convex, $c\in\R$, $d_1\ne d_2$, and
\[
f=\partial_t+c\,\partial_x-d_1,\qquad g=\partial_t+c\,\partial_x-d_2 .
\]
Let $\ell(s)$ be the length of the chord of $\Omega$ along the characteristic line $x-ct=s$, and $\mathcal H_i=\ker_{L^2}f_i$. The following are equivalent:
\begin{enumerate}[label=(\alph*)]
\item the chords do not shrink: $\inf_s\ell(s)>0$;
\item every $u\in L^2$ with $fgu=0$ has components in $L^2$, with
\[
\|u_1\|^2+\|u_2\|^2\le C\,\|u\|^2 ;
\]
\item the sum $\mathcal H_1+\mathcal H_2$ is closed in $L^2$;
\item the Friedrichs cosine of the two spaces is smaller than one, $c_F(\mathcal H_1,\mathcal H_2)<1$.
\end{enumerate}
\dimref{C12-T3}
\end{theorem}

\begin{example}{Rectangles}
On a rectangle $I\times J$ with $c\ne0$, the chords in direction $(1,c)$ shrink at two corners. So the global $L^2$ decomposition always fails ($c_F=1$). For $c=0$ the chords are vertical, all of length $|I|$, and the decomposition holds.
\end{example}

\begin{figure}[!htb]
\centering
\begin{tikzpicture}[scale=1.1]
\draw[->] (-0.3,0)--(3.6,0); \draw[->] (0,-0.3)--(0,2.8);
\draw[cA,very thick,->] (0,0)--(3.2,0.5) node[right,font=\small]{$M_1=\ker Q_1$};
\draw[cB,very thick,->] (0,0)--(1.9,2.4) node[above,font=\small]{$M_2=\ker Q_2$};
\draw[cG] (0.9,0.14) arc (8.9:51.6:0.9); \node[font=\small] at (1.12,0.55) {$\theta$};
\node[font=\small,anchor=west] at (4.3,1.6) {$c_F=\cos\theta$};
\node[font=\small,anchor=west] at (4.3,1.0) {$\gamma_2=1-c_F$};
\end{tikzpicture}
\caption{The Friedrichs angle $\theta$ between two component spaces. A larger angle means a better-conditioned recovery of the pieces from their sum. As $\theta\to0$ the recovery becomes unstable.}\label{fig:friedrichs}
\end{figure}

\section{Any number of components}

For $N$ components the right quantity is the conditioning of the sum map
\[
S:\ \bigoplus_{j=1}^N\ker_{H^m}Q_j\to H^m,\qquad S(u_1,\dots,u_N)=\sum_{j=1}^Nu_j .
\]

\begin{definition}{Conditioning number $\gamma_N$}\label{def:gamma}
Let $(\ker S)^\perp$ be the tuples orthogonal to the ambiguity of Theorem~\ref{thm:factor}. Then
\[
\gamma_N=\inf_{\substack{v\in(\ker S)^\perp\\ v\ne0}}\ \frac{\|Sv\|_{H^m}^2}{\|v\|^2},
\qquad \|v\|^2=\sum_j\|u_j\|_{H^m}^2 .
\]
Equivalently, $\gamma_N$ is the bottom of the spectrum of $S^*S$ on $(\ker S)^\perp$: the square of the smallest nonzero singular value of $S$.
\end{definition}

\begin{intuition}
$\gamma_N$ is the worst ratio between the size of the sum and the size of the pieces. If $\gamma_N>0$, the pieces cannot nearly cancel each other:
\[
\|v\|\le\gamma_N^{-1/2}\,\|Sv\| ,
\]
and recovering them from $u$ is stable. If $\gamma_N=0$, there are large pieces with a tiny sum, and the recovery is unstable. It is exactly the smallest singular value of a least-squares problem.
\end{intuition}

\begin{theorem}{Regularity criterion \led{C12-P7} \proved}\label{thm:hm}
\begin{enumerate}[label=(\alph*)]
\item Every $u\in\ker_{H^m}P$ has components in $H^m$ if and only if the sum $\sum_j\ker_{H^m}Q_j$ is \emph{closed} and \emph{dense} in $\ker_{H^m}P$.
\item The sum is closed if and only if $\gamma_N>0$. Then the minimal-energy split satisfies, optimally,
\[
\sum_j\|u_j\|_{H^m}^2\le\gamma_N^{-1}\,\|u\|_{H^m}^2 .
\]
\item For $N=2$, $\gamma_2=1-c_F$, the Friedrichs cosine of the two spaces after removing their intersection.
\item For $N\ge3$ the pairwise cosines do not decide. There are three closed subspaces of a Hilbert space with trivial pairwise intersections and closed pairwise sums ($c_F\le\tfrac12$), whose triple sum is not closed.
\end{enumerate}
\dimref{C12-P7}
\end{theorem}

The condition has two parts. Closedness is the conditioning above. Density is a separate question, and on convex regions it is always true.

\begin{theorem}{Density on convex regions \led{OPF5B-T1} \modulo{H\"ormander, Stein}}\label{thm:density}
Let $\Omega\subset\R^2$ be bounded and convex. For every operator $P$, every coprime factorisation and every $m\ge0$, the sum $\sum_j\ker_{H^m}Q_j$ is dense in $\ker_{H^m}P$.

\dimref{OPF5B-T1}
\end{theorem}

So on a convex region only $\gamma_N>0$ remains to be checked.

\begin{caution}[title=Added in the complete edition (Phase 10)]
Convexity cannot simply be dropped: on a region with a hole the sum of kernels is \emph{not} dense in general (Proposition~\ref{prop:s10ap1}, the Laplacian on an annulus), and the extension of this theorem to non-convex regions is false. For simply connected Lipschitz regions density was proved in Phase~10 modulo two cited facts (Chapter~\ref{ch:phase10}).
\end{caution}

\begin{example}{Cases where the $H^m$ decomposition holds \led{C12-P4} \proved}
\textbf{$P=\tau\xi$ on a rectangle.} The components are a function of $t$ and a function of $x$. After removing the constants their $H^m$ inner product vanishes, so $c_F=0$ and $\gamma_2=1$.

\textbf{The wave equation on a square.} In characteristic coordinates $s=x+t$, $r=x-t$ the square becomes a rhombus. It can be covered by five pieces on each of which one component is controlled by $u$.

\textbf{Heat and transport on a rectangle} \led{OPR2-T1} \pw. For $P=(\tau-\xi^2)(\tau+c\xi)$ the sum is dense in every $H^m$, so the decomposition holds exactly when $c_F<1$.

The same operator can behave well on one region and badly on another: regularity depends on the \emph{geometry of $\Omega$ against the characteristics}, not on $P$ alone.

\dimref{C12-P4}
\end{example}

\begin{proposition}{Real transports, $m=0$ \led{OPR2-T2} \pw}\label{prop:transports}
Let $\Omega$ be convex and bounded, and $Q_j=v_j\cdot\nabla-d_j$ real transports, grouped into classes of parallel directions. Assume that of any two classes at least one has chords bounded below. Then $\sum_j\ker_{L^2}Q_j$ is closed if and only if every class with at least two factors has chords bounded below. In particular, for such transports, closed pairwise sums imply a closed triple sum.

\dimref{OPR2-T2}
\end{proposition}

\begin{caution}[title={Three or more components: what is still open}]
Whether the abstract counterexample of Theorem~\ref{thm:hm}(d) can occur for \emph{differential} factors is open in general \openq{} \led{C12-Q2}. Proposition~\ref{prop:transports} rules it out for real transports in $L^2$.
\end{caution}

\begin{keypoint}
Algebraically the split always exists (Theorem~\ref{thm:factor}). Analytically the components are as regular as $u$ exactly when the component spaces are separated by a positive angle, $\gamma_N>0$; density is automatic on convex regions. Everything delicate happens at the boundary of $\Omega$ (shrinking chords) or in its topology (holes).
\end{keypoint}
